  
%
% $Id: slides.tex 4228 2006-06-21 21:55:12Z jjamor $
%
%
% Compilar a .pdf con LaTeX (pdflatex)
% Es necesario instalar Beamer (paquete latex-beamer en Debian)
%



\documentclass{beamer}
\usetheme{Warsaw}
\usebackgroundtemplate{\includegraphics[width=\paperwidth]{format/libresoft-bg.png}}
\usepackage[english]{babel}
\usepackage[latin1]{inputenc}
\usepackage{graphics}
\usepackage{amssymb} % Simbolos matematicos
\usepackage{url}
%\usepackage{html}
%\usepackage{hthtml}

%\definecolor{libresoftgreen}{RGB}{162,190,43}
%\definecolor{libresoftblue}{RGB}{0,98,143}

%\setbeamercolor{titlelike}{bg=libresoftgreen}

%% Metadatos del PDF.
\hypersetup{
  pdftitle={Presentacion},
  pdfauthor={Roberto Calvo Palomino},
  pdfcreator={GSyC/Libresoft},
  pdfproducer=PDFLaTeX,
  pdfsubject={nn},
}
%%


\AtBeginSection[]
{
  \begin{frame}<presentation>
    \frametitle{Index}
    \tableofcontents[current]
  \end{frame}
}


\begin{document}

\title{Curso de Desarrollo en Android }
\subtitle{Presentacion}
\institute{\\rocapal@libresoft.es\\
GSyC/Libresoft}
\author{Roberto Calvo Palomino}
\date{\today}

\frame{
\maketitle
\begin{center}
\includegraphics[width=6cm]{format/gsyc-urjc}
\end{center}
}


%\title{Titulo corto}
%\author{Autores abreviado}


%% LICENCIA DE REDISTRIBUCION DE LAS TRANSPAS
\frame{
~
\vspace{4cm}

\begin{flushright}
{\tiny
(cc) 2010 Roberto Calvo Palomino. \\
Some rights reserved. This document is distributed under the Creative \\
            Commons Attribution-ShareAlike 3.0 licence, available in \\
            http://creativecommons.org/licenses/by-sa/3.0/


}
\end{flushright}
}


\section{Introducci�n}

\begin{frame}
\frametitle{Informaci�n General}
\begin{itemize}
\item Objetivo: Aprender a programar en dispositivos m�viles
\item Inicio: 4 de Marzo 2011 - Fin: 6 de Mayo 2011
\item Horario: de 16:00 a 21:00 los viernes
\item Lugar: Aula 007, Laboratorios II - URJC (Fuenlabrada)
\item Material: Ordenador, m�vil y libro de referencia
\item Al finalizar, obtendr�is el certificado de aprovechamiento del
  curso.
\end{itemize}
\end{frame}

\section{Temario}
\begin{frame}
\frametitle{Clases}
\begin{itemize}
\item Introduccion a Android. Filosof�a
\item Instalaci�n y Configuraci�n del SDK
\item Entorno y anatom�a de una app Android
\item Layouts y Activity
\item Tipos de Aplicaciones
\item Mapas y Geolocalizacion
\item Gesti�n de datos locales y remotos
\item Audio y Video. Streaming
\item Trabajando en background
\item Integra en Facebook tu aplicaci�n
\item Como subir aplicaciones al Market
\end{itemize}

\end{frame}


\begin{frame}
\frametitle{Charlas}

\begin{itemize}
\item Fragmentaci�n en Android (LibreSoft)
\item Visi�n de una operadora: Movistar
\item Crea tu propia ROM de Android (�ndago)
\item Promociona tu aplicaci�n m�vil (Google)
\end{itemize}

\end{frame}


\section{Metodolog�a}
\begin{frame}
\frametitle{Clases}

\begin{itemize}
\item Clases muy pr�cticas. Desarrollo de c�digo
\item Libro de referencia
\item En casa: leer libro y algo de c�digo
\item En clase: teor�a, c�digo y comentar c�digo
\item Traspas, c�digo y foros en moodle
\end{itemize}

\begin{center}
\htmladdnormallink{Moodle:
  http://docencia.etsit.urjc.es/moodle/course/view.php?id=186 }{http://docencia.etsit.urjc.es/moodle/course/view.php?id=186}
\end{center}

\end{frame}


pregu
\begin{frame}
\frametitle{Equipos y Material}

\begin{itemize}
\item Puedes usar tu propio port�til
\item Puedes usar los ordenadores del laboratorio
\item WIFI para el entorno:
\begin{itemize}
\item ESSID: wireless-urjc  (auth web)
\item usuario: congreso@urjc.es
\item contrase�a: c0ngres0
\end{itemize}
\item HashTag Twitter: \#clau (Curso Libresoft Android Urjc)
\item Canal IRC: \#clau en irc.freenode.net
\end{itemize}

\end{frame}



\begin {frame}
\frametitle{Preguntas}

\center{\huge{PREGUNTAS/DUDAS}}

\end{frame}


\end{document}
